\documentclass[11pt,a4paper]{article}
\usepackage[T1]{fontenc}
\usepackage[utf8]{inputenc}
\usepackage{amsmath,amssymb,amsthm}
\usepackage[margin=1in]{geometry}
\usepackage[colorlinks=true,linkcolor=blue,urlcolor=blue]{hyperref}
\theoremstyle{plain}
\newtheorem{theorem}{Theorem}
\newtheorem{lemma}[theorem]{Lemma}
\newtheorem{proposition}[theorem]{Proposition}
\newtheorem{corollary}[theorem]{Corollary}
\theoremstyle{definition}
\newtheorem{remark}[theorem]{Remark}

\title{The empirical recurrence for OEIS A253373, proved by transfer matrix}
\author{Adrian Perez Fontelles\\ \small Independent researcher}
\date{9 September 2026}
\begin{document}
\maketitle

\begin{abstract}
OEIS A253373 counts the arrays over $\{0,\dots,2\}$ whose $2\times2$ subblocks carry a
statistic that is monotone along named directions of the subblock grid, and carries an
empirical recurrence of order $45$ contributed by R.~H.~Hardin. It is true, and it is
decidable rather than empirical. Every direction named moves the subblock row by at most one,
so a pair of consecutive array rows is a state, the count is a walk count on $S=192474$ of them,
and the conjectured recurrence is then settled by a finite exact computation.
\end{abstract}

\noindent\small 2020 Mathematics Subject Classification. 05A15, 05B45, 06A07.\normalsize

\section{The conjecture}

OEIS A253373 is ``Number of (n+1)X(6+1) 0..2 arrays with every 2X2 subblock sum nondecreasing horizontally, vertically and antidiagonally ne-to-sw.''. It has offset $1$ and begins
\[
192474, 5125790, 128407497, 943840336, 8802760516, 35861948220, 195248382407, 603872541953,\ \dots
\]
The entry states, as a conjecture contributed by its author and never marked settled:
\begin{quote}
Empirical: a(n) = 9*a(n-1) +18*a(n-2) -390*a(n-3) +363*a(n-4) +7557*a(n-5) -16924*a(n-6) -85080*a(n-7) +289038*a(n-8) +592626*a(n-9) -3022596*a(n-10) -2302692*a(n-11) +22010974*a(n-12) +263586*a(n-13) -118108608*a(n-14) +60625464*a(n-15) +480363315*a(n-16) -453795339*a(n-17) -1497838078*a(n-18) +2041953882*a(n-19) +3566209239*a(n-20) -6619257735*a(n-21) -6304597140*a(n-22) +16357594512*a(n-23) +7550240128*a(n-24) -31527338352*a(n-25) -3791317152*a(n-26) +47758481952*a(n-27) -6659648352*a(n-28) -56728278048*a(n-29) +20129722880*a(n-30) +52207707264*a(n-31) -28856183808*a(n-32) -36302235648*a(n-33) +27669476352*a(n-34) +18144427008*a(n-35) -18828638464*a(n-36) -5799153408*a(n-37) +9055982592*a(n-38) +719677440*a(n-39) -2944438272*a(n-40) +244629504*a(n-41) +582598656*a(n-42) -119439360*a(n-43) -53084160*a(n-44) +15925248*a(n-45) for n>67
\end{quote}
As of the ``Last modified'' line on the live entry (Jul 23 2025 14:00:54, revision 6) this is still
recorded as empirical, and nothing on the entry records it as proved.

\section{What is being counted}

The $2\times2$ subblocks of an $(n+1)\times7$ array over $\{0,\dots,2\}$ form a grid
of $n$ rows and $6$ columns; write $\sigma(i,j)$ for the value the entry's statistic takes on
the subblock in position $(i,j)$, a function of its four entries
$\begin{pmatrix}p&q\\r&s\end{pmatrix}$ alone. The condition is
\[
\sigma(i,j)\le\sigma(i,j+1),\qquad \sigma(i,j)\le\sigma(i+1,j),\qquad \sigma(i,j)\le\sigma(i+1,j-1)
\]
whenever both positions exist.

Three points of the English are fixed by the entries' own terms rather than guessed:
``ne-to-sw antidiagonally'' is the offset $(+1,-1)$ on the subblock grid, since the offset
$(+1,+1)$ gives $164$ where the corresponding entry publishes $126$; ``diagonal minus
antidiagonal sum'' is $(p+s)-(q+r)$; and ``ne-sw antidiagonal difference'' is $q-r$. Each was
checked by counting the small arrays directly and comparing with the entry.

\section{The count is a walk count}

A pair of consecutive array rows determines one whole row of the subblock grid. Take those
pairs as states. The comparisons with offset $(0,\pm1)$ live inside a single subblock row, so
they decide which pairs are admissible at all; the comparisons with offset $(\pm1,\cdot)$
relate two consecutive subblock rows, so they decide which pair may follow which. A step of the
walk appends one array row, turning the pair $(r,s)$ into $(s,t)$, and settles every comparison
between the two subblock rows at once. An $(n+1)$-row array is therefore a walk of $n-1$ steps
and every walk comes from exactly one array, so
\[
a(n)\;=\;\iota^{\!\top}M^{\,n-1}\tau,\qquad \iota=\tau=(1,\dots,1)^{\!\top},
\]
on the $S=192474$ admissible pairs. In particular $a$ is $C$-finite. The alignment of the walk
index with the entry's own $n$ was checked against its published terms rather than assumed.

\section{The proof}

\begin{lemma}\label{lem:crit}
Let $q(t)=t^{45}-\sum_i c_it^{45-i}$ be the characteristic polynomial of the
conjectured recurrence. The residual $a(n)-\sum_ic_ia(n-i)$ equals
$\iota^{\!\top}M^{\,j}q(M)\tau$ for the corresponding walk index $j$, so the recurrence
holds from some point on if and only if $\iota^{\!\top}M^{j}q(M)\tau=0$ for all large $j$,
and it is enough to examine $j<S$.
\end{lemma}

\begin{proof}
Factor $M^{j}$ out of $M^{j+45}-\sum_ic_iM^{j+45-i}$. For the bound, $M^{S}$
is an integer combination of $I,M,\dots,M^{S-1}$ by Cayley--Hamilton, so the numbers
$u_j=\iota^{\!\top}M^{j}q(M)\tau$ obey the monic recurrence given by the characteristic
polynomial of $M$; $S$ consecutive zeros force every later one, and the last nonzero $u_j$
pins the threshold exactly.
\end{proof}

\begin{theorem}
\[
a(n)\;=\;9\,a(n-1) + 18\,a(n-2) - 390\,a(n-3) + 363\,a(n-4) + 7557\,a(n-5) - 16924\,a(n-6) - 85080\,a(n-7) + 289038\,a(n-8) + 592626\,a(n-9) - 3022596\,a(n-10) - 2302692\,a(n-11) + 22010974\,a(n-12) + 263586\,a(n-13) - 118108608\,a(n-14) + 60625464\,a(n-15) + 480363315\,a(n-16) - 453795339\,a(n-17) - 1497838078\,a(n-18) + 2041953882\,a(n-19) + 3566209239\,a(n-20) - 6619257735\,a(n-21) - 6304597140\,a(n-22) + 16357594512\,a(n-23) + 7550240128\,a(n-24) - 31527338352\,a(n-25) - 3791317152\,a(n-26) + 47758481952\,a(n-27) - 6659648352\,a(n-28) - 56728278048\,a(n-29) + 20129722880\,a(n-30) + 52207707264\,a(n-31) - 28856183808\,a(n-32) - 36302235648\,a(n-33) + 27669476352\,a(n-34) + 18144427008\,a(n-35) - 18828638464\,a(n-36) - 5799153408\,a(n-37) + 9055982592\,a(n-38) + 719677440\,a(n-39) - 2944438272\,a(n-40) + 244629504\,a(n-41) + 582598656\,a(n-42) - 119439360\,a(n-43) - 53084160\,a(n-44) + 15925248\,a(n-45)
\]
for every $n>67$.
\end{theorem}

\begin{proof}
The vector $w=q(M)\tau$ was formed by $45$ matrix--vector products in exact integer
arithmetic and $\iota^{\!\top}M^{j}w$ evaluated for $j=0,1,\dots$, again exactly; those
integers vanish from the index corresponding to $n=67+1$ onwards and the run continued
until the working vector was identically zero, which settles every larger $j$ at once.
\end{proof}

\section{Verification}

Three checks, each independent of the algebra above.

First, the model reproduces all $16$ terms the entry publishes, exactly, in integer
arithmetic. This is what ties the model to the entry: the digraph is built by reading the
entry's English, and a misreading gives different counts.

Second, the conjectured recurrence was evaluated directly on those published terms, with no
matrices involved, and holds wherever the proved range and the data overlap.

Third, the annihilation test of Lemma~\ref{lem:crit} was carried out in exact integer
arithmetic on the whole vector rather than on a sampled prefix.

\begin{thebibliography}{9}
\bibitem{oeis} The OEIS Foundation, \emph{The On-Line Encyclopedia of Integer
Sequences}, \url{https://oeis.org}, 2026. Sequence A253373.
\bibitem{stanley} R.~P.~Stanley, \emph{Enumerative Combinatorics}, Volume 1, 2nd ed.,
Cambridge University Press, 2011. (Transfer-matrix method, Section 4.7.)
\bibitem{davey} B.~A.~Davey and H.~A.~Priestley, \emph{Introduction to Lattices and
Order}, 2nd ed., Cambridge University Press, 2002.
\bibitem{horn} R.~A.~Horn and C.~R.~Johnson, \emph{Matrix Analysis}, 2nd ed., Cambridge
University Press, 2013.
\end{thebibliography}

\end{document}
